%% toc: P2 13 13.1,13.2,13.3
\section{Proofs for the derivation maps between the classical calculi}
\label{sec:annex-p02-w02b-p02-010-translations}

Notation is that of \cref{def:p02-w02b-p02-010-setting}.  In Hilbert
derivations below, a line labelled by an item of
\cref{thm:p02-w02b-p02-010-hilbert-theorems} or
\cref{thm:p02-w02b-p02-010-rule-images}, or by an earlier displayed
derivation, stands for that derivation inserted in place; by
\cref{thm:p02-w02b-p02-010-hilbert-structural}(1)--(2) the result is a
Hilbert derivation from the stated premises.  ``DT'' cites the deduction
theorem, \cref{thm:p02-hilbert-deduction}, proved in \cref{annex:p2.calc.2};
its right-to-left half is one MP step.  Natural-deduction derivations are
written as numbered lines in which each line names the lines of its
premises and each discharge names its label.

\begin{annexitem}{annex:p02-w02b-p02-010-hilbert-structural}{thm:p02-w02b-p02-010-hilbert-structural}{Structural properties of Hilbert derivability}
\begin{proof}
(1) A derivation from \(\Gamma\) is a finite sequence each of whose lines is
in \(\Gamma\), is an axiom instance, or follows from two earlier lines by MP.
Each such line satisfies the same condition with \(\Delta\supseteq\Gamma\) in
place of \(\Gamma\), so the same sequence is a derivation from \(\Delta\).

(2) Let \(\theta_0,\ldots,\theta_n=\theta\) be a derivation from
\(\Gamma\cup\Delta\).  For each \(i\) with \(\theta_i\in\Delta\setminus\Gamma\)
choose a derivation \(\rho_i\) of \(\theta_i\) from \(\Gamma\); its last line
is \(\theta_i\).  Replace each such line \(\theta_i\) by the segment
\(\rho_i\).  A line inside an inserted segment is justified within that segment,
since an MP line of \(\rho_i\) cites two earlier lines of \(\rho_i\).  A line
\(\theta_m\) of the original sequence that was not replaced is in
\(\Gamma\), is an axiom instance, or follows by MP from earlier
\(\theta_j,\theta_k\); each of \(\theta_j,\theta_k\) is still present before
\(\theta_m\), either as itself or as the last line of its segment.  The last
line of the new sequence is \(\theta\), either unchanged or as the last
line of \(\rho_n\).  Hence the new sequence is a derivation of \(\theta\)
from \(\Gamma\).

(3) A derivation has finitely many lines.  Its lines that are cited as
members of \(\Gamma\) form a finite set \(\Gamma_0\subseteq\Gamma\), and
every other line is an axiom instance or an MP line, so the same sequence
is a derivation from \(\Gamma_0\).
\end{proof}
\end{annexitem}

\begin{annexitem}{annex:p02-w02b-p02-010-hilbert-theorems}{thm:p02-w02b-p02-010-hilbert-theorems}{Derived Hilbert theorems}
\begin{proof}
Each item gives a derivation; schema instances are named by the sheet's
labels A1, A2, A3, where A3 is
\((\neg\psi\to\neg\varphi)\to((\neg\psi\to\varphi)\to\psi)\).

(1) This is the derivation in \cref{annex:p2.calc.2}:
\[
\begin{array}{rll}
1.&\varphi\to((\varphi\to\varphi)\to\varphi)&\text{A1}\\
2.&(\varphi\to((\varphi\to\varphi)\to\varphi))\to
  ((\varphi\to(\varphi\to\varphi))\to(\varphi\to\varphi))&\text{A2}\\
3.&(\varphi\to(\varphi\to\varphi))\to(\varphi\to\varphi)&\text{MP }1,2\\
4.&\varphi\to(\varphi\to\varphi)&\text{A1}\\
5.&\varphi\to\varphi&\text{MP }4,3
\end{array}
\]

(2) From \(\{\varphi\to\psi,\psi\to\chi,\varphi\}\) the lines
\(\varphi\), \(\varphi\to\psi\), \(\psi\) (MP), \(\psi\to\chi\), \(\chi\) (MP)
form a derivation of \(\chi\).  DT, discharging \(\varphi\), gives
\(\{\varphi\to\psi,\psi\to\chi\}\derives_{\mathsf H_{\mathsf{CPL}}}\varphi\to\chi\).

(3) From \(\{\neg\neg\varphi\}\):
\[
\begin{array}{rll}
1.&\neg\neg\varphi&\text{premise}\\
2.&\neg\neg\varphi\to(\neg\varphi\to\neg\neg\varphi)&\text{A1}\\
3.&\neg\varphi\to\neg\neg\varphi&\text{MP }1,2\\
4.&(\neg\varphi\to\neg\neg\varphi)\to((\neg\varphi\to\neg\varphi)\to\varphi)
  &\text{A3 with }\psi\mapsto\varphi,\ \varphi\mapsto\neg\varphi\\
5.&(\neg\varphi\to\neg\varphi)\to\varphi&\text{MP }3,4\\
6.&\neg\varphi\to\neg\varphi&\text{item 1}\\
7.&\varphi&\text{MP }6,5
\end{array}
\]
DT gives \(\derives_{\mathsf H_{\mathsf{CPL}}}\neg\neg\varphi\to\varphi\).

(4) From \(\{\varphi\}\):
\[
\begin{array}{rll}
1.&\neg\neg\neg\varphi\to\neg\varphi&\text{item 3 for }\neg\varphi\\
2.&(\neg\neg\neg\varphi\to\neg\varphi)\to((\neg\neg\neg\varphi\to\varphi)\to\neg\neg\varphi)
  &\text{A3 with }\psi\mapsto\neg\neg\varphi\\
3.&(\neg\neg\neg\varphi\to\varphi)\to\neg\neg\varphi&\text{MP }1,2\\
4.&\varphi&\text{premise}\\
5.&\varphi\to(\neg\neg\neg\varphi\to\varphi)&\text{A1}\\
6.&\neg\neg\neg\varphi\to\varphi&\text{MP }4,5\\
7.&\neg\neg\varphi&\text{MP }6,3
\end{array}
\]
DT gives \(\derives_{\mathsf H_{\mathsf{CPL}}}\varphi\to\neg\neg\varphi\).

(5) From \(\{\neg\varphi,\varphi\}\):
\[
\begin{array}{rll}
1.&\neg\varphi&\text{premise}\\
2.&\neg\varphi\to(\neg\psi\to\neg\varphi)&\text{A1}\\
3.&\neg\psi\to\neg\varphi&\text{MP }1,2\\
4.&\varphi&\text{premise}\\
5.&\varphi\to(\neg\psi\to\varphi)&\text{A1}\\
6.&\neg\psi\to\varphi&\text{MP }4,5\\
7.&(\neg\psi\to\neg\varphi)\to((\neg\psi\to\varphi)\to\psi)&\text{A3}\\
8.&(\neg\psi\to\varphi)\to\psi&\text{MP }3,7\\
9.&\psi&\text{MP }6,8
\end{array}
\]
DT discharging \(\varphi\) and then \(\neg\varphi\) gives
\(\derives_{\mathsf H_{\mathsf{CPL}}}\neg\varphi\to(\varphi\to\psi)\).

(6) From \(\{\neg\psi\to\neg\varphi,\varphi\}\):
\[
\begin{array}{rll}
1.&\neg\psi\to\neg\varphi&\text{premise}\\
2.&(\neg\psi\to\neg\varphi)\to((\neg\psi\to\varphi)\to\psi)&\text{A3}\\
3.&(\neg\psi\to\varphi)\to\psi&\text{MP }1,2\\
4.&\varphi&\text{premise}\\
5.&\varphi\to(\neg\psi\to\varphi)&\text{A1}\\
6.&\neg\psi\to\varphi&\text{MP }4,5\\
7.&\psi&\text{MP }6,3
\end{array}
\]
DT twice gives
\(\derives_{\mathsf H_{\mathsf{CPL}}}(\neg\psi\to\neg\varphi)\to(\varphi\to\psi)\).

(7) From \(\{\varphi\to\psi\}\):
\[
\begin{array}{rll}
1.&\neg\neg\varphi\to\varphi&\text{item 3}\\
2.&\varphi\to\psi&\text{premise}\\
3.&\neg\neg\varphi\to\psi&\text{item 2 from }1,2\\
4.&\psi\to\neg\neg\psi&\text{item 4}\\
5.&\neg\neg\varphi\to\neg\neg\psi&\text{item 2 from }3,4\\
6.&(\neg\neg\varphi\to\neg\neg\psi)\to(\neg\psi\to\neg\varphi)
  &\text{item 6 with }\psi\mapsto\neg\varphi,\ \varphi\mapsto\neg\psi\\
7.&\neg\psi\to\neg\varphi&\text{MP }5,6
\end{array}
\]
Lines 3 and 5 use item 2 through
\cref{thm:p02-w02b-p02-010-hilbert-structural}(2), with the premises of
item 2 derived at the cited earlier lines.  DT gives
\(\derives_{\mathsf H_{\mathsf{CPL}}}(\varphi\to\psi)\to(\neg\psi\to\neg\varphi)\).

(8) From \(\{\varphi\to\psi,\neg\varphi\to\psi\}\):
\[
\begin{array}{rll}
1.&\varphi\to\psi&\text{premise}\\
2.&(\varphi\to\psi)\to(\neg\psi\to\neg\varphi)&\text{item 7}\\
3.&\neg\psi\to\neg\varphi&\text{MP }1,2\\
4.&\neg\varphi\to\psi&\text{premise}\\
5.&(\neg\varphi\to\psi)\to(\neg\psi\to\neg\neg\varphi)&\text{item 7 for }\neg\varphi\\
6.&\neg\psi\to\neg\neg\varphi&\text{MP }4,5\\
7.&(\neg\psi\to\neg\neg\varphi)\to((\neg\psi\to\neg\varphi)\to\psi)
  &\text{A3 with }\varphi\mapsto\neg\varphi\\
8.&(\neg\psi\to\neg\varphi)\to\psi&\text{MP }6,7\\
9.&\psi&\text{MP }3,8
\end{array}
\]
DT twice gives
\(\derives_{\mathsf H_{\mathsf{CPL}}}(\varphi\to\psi)\to((\neg\varphi\to\psi)\to\psi)\).
No item uses a later item, so the construction is well founded.
\end{proof}
\end{annexitem}

\begin{annexitem}{annex:p02-w02b-p02-010-rule-images}{thm:p02-w02b-p02-010-rule-images}{Hilbert images of the natural-deduction rules}
\begin{proof}
Items of \cref{thm:p02-w02b-p02-010-hilbert-theorems} are cited by number.

(i) From \(\{A,A\to\neg B\}\) one MP step gives \(\neg B\), so DT gives
\(\{A\}\vdash(A\to\neg B)\to\neg B\).  Item 7 with
\(\varphi\mapsto A\to\neg B\), \(\psi\mapsto\neg B\) is
\(((A\to\neg B)\to\neg B)\to(\neg\neg B\to\neg(A\to\neg B))\); MP gives
\(\neg\neg B\to\neg(A\to\neg B)\).  From the premise \(B\), item 4 and MP
give \(\neg\neg B\), and a final MP gives \(\neg(A\to\neg B)\).

(ii) Item 5 gives \(\neg A\to(A\to\neg B)\).  Item 7 with
\(\varphi\mapsto\neg A\), \(\psi\mapsto A\to\neg B\) and MP give
\(\neg(A\to\neg B)\to\neg\neg A\).  MP with the premise
\(\neg(A\to\neg B)\) gives \(\neg\neg A\), and item 3 with MP gives \(A\).
For the second conjunct, A1 gives \(\neg B\to(A\to\neg B)\); item 7 with
\(\varphi\mapsto\neg B\), \(\psi\mapsto A\to\neg B\) and MP give
\(\neg(A\to\neg B)\to\neg\neg B\); MP with the premise gives
\(\neg\neg B\), and item 3 with MP gives \(B\).

(iii) From \(\{A,\neg A\}\), item 5 \(\neg A\to(A\to B)\) and two MP steps
give \(B\); DT discharging \(\neg A\) gives \(\{A\}\vdash\neg A\to B\).
From \(\{B\}\), the A1 instance \(B\to(\neg A\to B)\) and MP give
\(\neg A\to B\).

(iv) Item 2 applied to \(\neg A\to B\) and \(B\to C\) gives
\(\neg A\to C\).  Item 8 with \(\varphi\mapsto A\), \(\psi\mapsto C\) is
\((A\to C)\to((\neg A\to C)\to C)\); MP with \(A\to C\) and then with
\(\neg A\to C\) gives \(C\).

(v) Item 7 with \(\varphi\mapsto A\), \(\psi\mapsto\top_{\mathsf H}\) is
\((A\to\neg\top_{\mathsf H})\to(\neg\neg\top_{\mathsf H}\to\neg A)\); MP with
the premise \(A\to\bot_{\mathsf H}\), which is \(A\to\neg\top_{\mathsf H}\),
gives \(\neg\neg\top_{\mathsf H}\to\neg A\).  Item 1 for \(p_0\) gives
\(\top_{\mathsf H}\), item 4 gives
\(\top_{\mathsf H}\to\neg\neg\top_{\mathsf H}\), and two MP steps give
\(\neg A\).

(vi) Item 5 with \(\varphi\mapsto A\), \(\psi\mapsto\bot_{\mathsf H}\) is
\(\neg A\to(A\to\bot_{\mathsf H})\); MP with \(\neg A\) and then with \(A\)
gives \(\bot_{\mathsf H}\).

(vii) Item 5 with \(\varphi\mapsto\top_{\mathsf H}\), \(\psi\mapsto C\) is
\(\neg\top_{\mathsf H}\to(\top_{\mathsf H}\to C)\); MP with the premise
\(\bot_{\mathsf H}=\neg\top_{\mathsf H}\) gives \(\top_{\mathsf H}\to C\), and
MP with \(\top_{\mathsf H}\) (item 1) gives \(C\).

(viii) The A3 instance with \(\psi\mapsto A\), \(\varphi\mapsto\top_{\mathsf H}\)
is \((\neg A\to\neg\top_{\mathsf H})\to((\neg A\to\top_{\mathsf H})\to A)\).
MP with the premise gives \((\neg A\to\top_{\mathsf H})\to A\).  From
\(\top_{\mathsf H}\) (item 1) and the A1 instance
\(\top_{\mathsf H}\to(\neg A\to\top_{\mathsf H})\), MP gives
\(\neg A\to\top_{\mathsf H}\), and a last MP gives \(A\).
\end{proof}
\end{annexitem}

\begin{annexitem}{annex:p02-w02b-p02-010-nd-structural}{thm:p02-w02b-p02-010-nd-structural}{Padding, unrestricted discharge and composition in natural deduction}
\begin{proof}
(1) Let \(D'\) consist of \(D\) and a new one-node assumption \(\alpha\),
joined by \(\land I\) to \(\chi\land\alpha\), followed by \(\land E_1\) to
\(\chi\).  Neither rule discharges anything, so the new leaf is open, every
discharge inside \(D\) is unchanged, and
\(\operatorname{oa}(D')=\operatorname{oa}(D)\cup\{\alpha\}\).

(2) For \(\to I\): let \(D\) witness
\(\Gamma\cup\{\alpha\}\derives_{\mathsf{ND}_{\mathsf{CPL}}}\beta\).  If \(D\)
has no open occurrence of \(\alpha\), replace it by the derivation \(D'\) of
part (1); this keeps the root and adds only \(\alpha\) to the open
assumptions.  Now \(D\) has an open occurrence of \(\alpha\).  Mark all its
open occurrences of \(\alpha\) with a label fresh from every label of \(D\),
by merging relabeling, and apply \(\to I\) with that label.  The discharge is
nonvacuous, and the open assumptions of the result are
\(\operatorname{oa}(D)\setminus\{\alpha\}\subseteq\Gamma\).  The cases
\(\neg I\) (root \(\bottom\), conclusion \(\neg\alpha\)) and RAA (discharged
formula \(\neg\alpha\), conclusion \(\alpha\)) are the same construction with
the named rule in place of \(\to I\) and, for RAA, \(\neg\alpha\) in place of
\(\alpha\) throughout.  For \(\lor E\): let \(D_1,D_2,D_3\) witness the three
premises.  Pad \(D_2\) with \(\alpha\) and \(D_3\) with \(\beta\) where
needed, rename the three derivations into pairwise disjoint label sets,
mark all open \(\alpha\) in \(D_2\) with a fresh label \(i\) and all open
\(\beta\) in \(D_3\) with a different fresh label \(j\), and apply \(\lor E\)
with \(i,j\).  Both discharges are nonvacuous, the marks touch only
occurrences in the intended premise derivation, and the open assumptions
of the result are contained in
\(\operatorname{oa}(D_1)\cup(\operatorname{oa}(D_2)\setminus\{\alpha\})
\cup(\operatorname{oa}(D_3)\setminus\{\beta\})\subseteq\Gamma\cup\Delta\cup\Sigma\).

(3) Let \(D\) witness \(\Gamma\cup\Sigma\derives_{\mathsf{ND}_{\mathsf{CPL}}}\chi\).
The finite set \(\operatorname{oa}(D)\setminus(\Gamma\cup\Delta)\) is
contained in \(\Sigma\); list it as \(\delta_1,\ldots,\delta_m\) and choose
derivations \(E_k\) of \(\delta_k\) with
\(\operatorname{oa}(E_k)\subseteq\Delta\).  Put \(D_0=D\), and obtain \(D_k\)
from \(D_{k-1}\) by \cref{lem:nd-substitution}, designating all open
occurrences of \(\delta_k\) and grafting copies of \(E_k\).  By that lemma,
\(\operatorname{oa}(D_k)\subseteq(\operatorname{oa}(D_{k-1})\setminus\{\delta_k\})
\cup\operatorname{oa}(E_k)\), so by induction on \(k\),
\(\operatorname{oa}(D_k)\subseteq\Gamma\cup\Delta\cup\{\delta_{k+1},\ldots,\delta_m\}\).
Discharge is by labels, as in \cref{def:p02-w02b-p02-010-setting}(2), and
each discharging application keeps its discharged occurrences.  Let \(r\)
be an application of \(D_{k-1}\) carrying label \(i\) for the formula
\(\gamma\).  Before grafting it discharges the occurrences of
\(\gamma\) marked \(i\) above the relevant premise; none of them is
replaced, because a replaced occurrence is open in \(D_{k-1}\) and so is
discharged by no application on its path.  After grafting, every leaf
inside a copy of \(E_k\) carries either no label or a label of the copy's
reserved set, which differs from \(i\); so \(r\) discharges no occurrence
inside a copy, and it discharges exactly the occurrences it discharged
before, which form a nonempty set.  Every application inside a copy of
\(E_k\) discharges the renamed images of the occurrences it discharged in
\(E_k\), by the renaming convention of the lemma.  So every discharge
remains nonvacuous, and \(D_m\) witnesses
\(\Gamma\cup\Delta\derives_{\mathsf{ND}_{\mathsf{CPL}}}\chi\).
\end{proof}
\end{annexitem}

\begin{annexitem}{annex:p02-w02b-p02-010-nd-translation}{thm:p02-w02b-p02-010-nd-translation}{Each formula is natural-deduction equivalent to its Hilbert translation}
\begin{proof}
Induct on \(\alpha\).  Write \(\vdash\) for
\(\derives_{\mathsf{ND}_{\mathsf{CPL}}}\), \(A=\operatorname{tr}_{\mathsf H}(\alpha)\),
\(B=\operatorname{tr}_{\mathsf H}(\beta)\), and use the induction hypotheses
\(\{\alpha\}\vdash A\), \(\{A\}\vdash\alpha\), \(\{\beta\}\vdash B\),
\(\{B\}\vdash\beta\) for the immediate subformulas.  ``(2)'' and ``(3)''
cite \cref{thm:p02-w02b-p02-010-nd-structural}; a set on the left of
\(\vdash\) may always be enlarged, since derivability only bounds the open
assumptions.

\emph{Variable \(p\).}  \(\operatorname{tr}_{\mathsf H}(p)=p\), and the
one-node derivation \(p\) gives both directions.

\emph{\(\bottom\).}  The one-node derivation \(\bottom\) followed by
\(\bottom E\) gives \(\{\bottom\}\vdash\neg(p_0\to p_0)\).  Conversely, the
assumption \([p_0]^1\) followed by \(\to I\) with label \(1\) is a closed
derivation of \(p_0\to p_0\); \(\neg E\) with the assumption
\(\neg(p_0\to p_0)\) gives \(\{\neg(p_0\to p_0)\}\vdash\bottom\).

\emph{\(\neg\alpha\).}  From \(\{A\}\vdash\alpha\) and the assumption
\(\neg\alpha\), \(\neg E\) gives \(\{A,\neg\alpha\}\vdash\bottom\); (2) for
\(\neg I\) discharging \(A\) gives \(\{\neg\alpha\}\vdash\neg A\).  From
\(\{\alpha\}\vdash A\) and the assumption \(\neg A\), \(\neg E\) gives
\(\{\alpha,\neg A\}\vdash\bottom\); (2) gives \(\{\neg A\}\vdash\neg\alpha\).

\emph{\(\alpha\to\beta\).}  From \(\{A\}\vdash\alpha\) and the assumption
\(\alpha\to\beta\), \(\to E\) gives \(\{A,\alpha\to\beta\}\vdash\beta\);
(3) with \(\{\beta\}\vdash B\) gives \(\{A,\alpha\to\beta\}\vdash B\); (2)
for \(\to I\) gives \(\{\alpha\to\beta\}\vdash A\to B\).  From
\(\{\alpha\}\vdash A\) and the assumption \(A\to B\), \(\to E\) gives
\(\{\alpha,A\to B\}\vdash B\); (3) with \(\{B\}\vdash\beta\) gives
\(\{\alpha,A\to B\}\vdash\beta\); (2) gives
\(\{A\to B\}\vdash\alpha\to\beta\).

\emph{\(\alpha\land\beta\), with translation \(N=\neg(A\to\neg B)\).}
\(\land E_1\) and \(\land E_2\) give \(\{\alpha\land\beta\}\vdash\alpha\) and
\(\vdash\beta\); by (3) and the hypotheses, \(\{\alpha\land\beta\}\vdash A\)
and \(\{\alpha\land\beta\}\vdash B\).  With the assumption \(A\to\neg B\),
\(\to E\) gives \(\neg B\), and \(\neg E\) with \(B\) gives
\(\{\alpha\land\beta,A\to\neg B\}\vdash\bottom\); (2) for \(\neg I\) gives
\(\{\alpha\land\beta\}\vdash N\).  Conversely, from the assumptions \(A\)
and \(\neg A\), \(\neg E\) and \(\bottom E\) give
\(\{A,\neg A\}\vdash\neg B\); (2) for \(\to I\) gives
\(\{\neg A\}\vdash A\to\neg B\); \(\neg E\) with the assumption \(N\) gives
\(\{N,\neg A\}\vdash\bottom\); (2) for RAA gives \(\{N\}\vdash A\), and (3)
gives \(\{N\}\vdash\alpha\).  From the one-node \(\neg B\), (2) for
\(\to I\) gives \(\{\neg B\}\vdash A\to\neg B\) even though \(A\) is unused;
\(\neg E\) with \(N\) gives \(\{N,\neg B\}\vdash\bottom\), RAA through (2)
gives \(\{N\}\vdash B\), and (3) gives \(\{N\}\vdash\beta\).  Finally
\(\land I\) gives \(\{N\}\vdash\alpha\land\beta\).

\emph{\(\alpha\lor\beta\), with translation \(M=\neg A\to B\).}  From
\(\{\alpha\}\vdash A\) and the assumption \(\neg A\), \(\neg E\) and
\(\bottom E\) give \(\{\alpha,\neg A\}\vdash B\); and \(\{\beta\}\vdash B\).
With the assumption \(\alpha\lor\beta\), (2) for \(\lor E\) gives
\(\{\alpha\lor\beta,\neg A\}\vdash B\), and (2) for \(\to I\) gives
\(\{\alpha\lor\beta\}\vdash M\).  Conversely, put
\(R=\neg(\alpha\lor\beta)\).  From \(\{A\}\vdash\alpha\), \(\lor I_1\) and
\(\neg E\) with the assumption \(R\) give \(\{A,R\}\vdash\bottom\), and (2)
for \(\neg I\) gives \(\{R\}\vdash\neg A\).  With the assumption \(M\),
\(\to E\) gives \(\{M,R\}\vdash B\); (3) with \(\{B\}\vdash\beta\), then
\(\lor I_2\) and \(\neg E\) with \(R\), give \(\{M,R\}\vdash\bottom\); and (2)
for RAA discharging \(R\) gives \(\{M\}\vdash\alpha\lor\beta\).

These clauses exhaust the formation rules of \(\mathsf{Fm}_{\mathsf{CPL}}\).
\end{proof}
\end{annexitem}

\begin{annexitem}{annex:p02-w02b-p02-010-hilbert-to-nd}{thm:p02-w02b-p02-010-hilbert-to-nd}{From Hilbert derivations to natural deduction}
\begin{proof}
By course-of-values induction on \(i\le n\) construct \(D_i\) with root
\(\theta_i\) and \(\operatorname{oa}(D_i)\subseteq\Gamma\).

\emph{Premise line} \(\theta_i\in\Gamma\): \(D_i\) is the one-node
assumption \(\theta_i\).

Discharge is by labels, as in \cref{def:p02-w02b-p02-010-setting}(2): a
line ``\(\to I,k\)'' or ``\(\mathsf{RAA},k\)'' discharges exactly the open
occurrences of its bracketed formula marked \(k\) above its premise.

\emph{A1 line} \(\varphi\to(\psi\to\varphi)\), for any
\(\varphi,\psi\), equal or not:
\[
\begin{array}{rll}
1.&[\varphi]^1&\text{assumption}\\
2.&[\psi]^2&\text{assumption}\\
3.&\varphi\land\psi&\land I,1,2\\
4.&\varphi&\land E_1,3\\
5.&\psi\to\varphi&\to I,2\\
6.&\varphi\to(\psi\to\varphi)&\to I,1
\end{array}
\]
Line 5 carries label \(2\) and discharges exactly the occurrence of
\(\psi\) marked \(2\), which is line 2; line 1 is marked \(1\), so it is not
discharged at line 5 even when \(\varphi=\psi\), and it is still open at
line 6, which discharges it.  Both labels mark an open occurrence of the
discharged formula, so neither discharge is vacuous, and the derivation is
closed.

\emph{A2 line}, with \(u=\varphi\to(\psi\to\chi)\) and
\(v=\varphi\to\psi\):
\[
\begin{array}{rll}
1.&[u]^1&\text{assumption}\\
2.&[\varphi]^3&\text{assumption}\\
3.&\psi\to\chi&\to E,1,2\\
4.&[v]^2&\text{assumption}\\
5.&[\varphi]^3&\text{assumption}\\
6.&\psi&\to E,4,5\\
7.&\chi&\to E,3,6\\
8.&\varphi\to\chi&\to I,3\\
9.&v\to(\varphi\to\chi)&\to I,2\\
10.&u\to(v\to(\varphi\to\chi))&\to I,1
\end{array}
\]
Line 8 carries label \(3\) and discharges exactly the occurrences of
\(\varphi\) marked \(3\), lines 2 and 5; line 9 discharges exactly the
occurrence of \(v\) marked \(2\), line 4; line 10 discharges exactly the
occurrence of \(u\) marked \(1\), line 1.  Each label marks an open
occurrence of the formula it discharges, so no discharge is vacuous, and
every leaf is discharged.

\emph{A3 line}, with \(u=\neg\psi\to\neg\varphi\) and
\(v=\neg\psi\to\varphi\):
\[
\begin{array}{rll}
1.&[u]^1&\text{assumption}\\
2.&[\neg\psi]^3&\text{assumption}\\
3.&\neg\varphi&\to E,1,2\\
4.&[v]^2&\text{assumption}\\
5.&[\neg\psi]^3&\text{assumption}\\
6.&\varphi&\to E,4,5\\
7.&\bottom&\neg E,6,3\\
8.&\psi&\mathsf{RAA},3\\
9.&v\to\psi&\to I,2\\
10.&u\to(v\to\psi)&\to I,1
\end{array}
\]
Line 8 carries label \(3\) and discharges exactly the occurrences of
\(\neg\psi\) marked \(3\), lines 2 and 5; line 9 discharges exactly line 4,
the occurrence of \(v\) marked \(2\); line 10 discharges exactly line 1, the
occurrence of \(u\) marked \(1\).  Each discharge is nonvacuous, and every
leaf is discharged.

\emph{MP line}: \(\theta_i\) follows from \(\theta_j\) and
\(\theta_k=\theta_j\to\theta_i\) with \(j,k<i\).  Rename \(D_k\) and \(D_j\)
into disjoint label sets and apply \(\to E\) with premises \(D_k\) and
\(D_j\).  The rule discharges nothing, so
\(\operatorname{oa}(D_i)=\operatorname{oa}(D_k)\cup\operatorname{oa}(D_j)\subseteq\Gamma\).

The axiom derivations are closed, so \(\operatorname{oa}(D_i)\subseteq\Gamma\)
in every case.  Taking \(i=n\) proves the consequence.  The rules used are
exactly those listed in the theorem: assumption, \(\land I\) and
\(\land E_1\) (A1), \(\to I\) (A1, A2, A3), \(\to E\) (A2, A3, MP),
\(\neg E\) and RAA (A3).
\end{proof}
\end{annexitem}

\begin{annexitem}{annex:p02-w02b-p02-010-nd-to-hilbert}{thm:p02-w02b-p02-010-nd-to-hilbert}{From natural deduction to Hilbert derivations}
\begin{proof}
Write \(\vdash\) for \(\derives_{\mathsf H_{\mathsf{CPL}}}\) and
\(\overline X\) for \(\operatorname{tr}_{\mathsf H}[X]\).  By induction on
the height of \(D\) we prove
\(\overline{\operatorname{oa}(D)}\vdash\operatorname{tr}_{\mathsf H}(\rho)\)
for the root \(\rho\) of \(D\).  Let \(D_1,D_2,D_3\) be the immediate
subderivations and \(\Gamma_k=\operatorname{oa}(D_k)\); capital letters
denote translations, \(A=\operatorname{tr}_{\mathsf H}(\alpha)\),
\(B=\operatorname{tr}_{\mathsf H}(\beta)\),
\(C=\operatorname{tr}_{\mathsf H}(\chi)\).  In each case the induction
hypotheses are brought to the common premise set
\(\overline{\operatorname{oa}(D)}\) by
\cref{thm:p02-w02b-p02-010-hilbert-structural}(1), and the cited item of
\cref{thm:p02-w02b-p02-010-rule-images} is attached by
\cref{thm:p02-w02b-p02-010-hilbert-structural}(2).

The open assumptions are accounted for by inclusions, because the sheet
discharges by labels (\cref{def:p02-w02b-p02-010-setting}(2)).  Let the last
application \(r\) of \(D\) have immediate subderivation \(D_k\).  An open
occurrence in \(D_k\) has, on its path to the root of \(D\), only the
applications on its path in \(D_k\), none of which discharges it, and \(r\)
itself.  If \(r\) does not discharge it, it is open in \(D\).  If \(r\)
discharges it, then \(r\) is a discharging application, \(D_k\) is the
derivation of a premise in which \(r\) discharges, and the occurrence is of
the formula \(r\) discharges there and carries \(r\)'s label for it.  Hence
\[
  \Gamma_k\subseteq\operatorname{oa}(D)
  \quad\text{if \(r\) discharges nothing in \(D_k\)},
  \qquad
  \Gamma_k\subseteq\operatorname{oa}(D)\cup\{\delta\}
  \quad\text{if \(r\) discharges occurrences of \(\delta\) in \(D_k\)}.
  \tag{\(\ast\)}
\]
No equation is claimed: when \(r\) leaves an occurrence of \(\delta\) in
\(D_k\) open, for instance an unmarked one, \(\delta\in\operatorname{oa}(D)\),
and \(\operatorname{oa}(D)\) may differ from \(\Gamma_k\setminus\{\delta\}\).
By (\(\ast\)) and \cref{thm:p02-w02b-p02-010-hilbert-structural}(1), an
induction hypothesis \(\overline{\Gamma_k}\vdash\theta\) yields
\(\overline{\operatorname{oa}(D)}\vdash\theta\) in the first case and
\(\overline{\operatorname{oa}(D)}\cup\{\operatorname{tr}_{\mathsf H}(\delta)\}\vdash\theta\)
in the second.  In the second case DT then gives
\(\overline{\operatorname{oa}(D)}\vdash\operatorname{tr}_{\mathsf H}(\delta)\to\theta\);
DT holds for every premise set, so it applies whether or not
\(\operatorname{tr}_{\mathsf H}(\delta)\) already belongs to
\(\overline{\operatorname{oa}(D)}\).  Write \(\Theta=\overline{\operatorname{oa}(D)}\).

\emph{Assumption} \(\alpha\): the one-line derivation of \(A\) from \(\{A\}\).

\emph{\(\land I\)}, \(\operatorname{oa}(D)=\Gamma_1\cup\Gamma_2\): the
hypotheses give \(A\) and \(B\); item (i) gives
\(\neg(A\to\neg B)=\operatorname{tr}_{\mathsf H}(\alpha\land\beta)\).

\emph{\(\land E_1\), \(\land E_2\)}: the hypothesis gives
\(\neg(A\to\neg B)\); item (ii) gives \(A\), respectively \(B\).

\emph{\(\lor I_1\), \(\lor I_2\)}: the hypothesis gives \(A\), respectively
\(B\); item (iii) gives \(\neg A\to B=\operatorname{tr}_{\mathsf H}(\alpha\lor\beta)\).

\emph{\(\lor E\)}, with major premise \(\alpha\lor\beta\) and discharges of
\(\alpha\) in \(D_2\) and of \(\beta\) in \(D_3\): by (\(\ast\)),
\(\Gamma_1\subseteq\operatorname{oa}(D)\),
\(\Gamma_2\subseteq\operatorname{oa}(D)\cup\{\alpha\}\) and
\(\Gamma_3\subseteq\operatorname{oa}(D)\cup\{\beta\}\).  The hypotheses
therefore give \(\Theta\vdash\neg A\to B\), \(\Theta\cup\{A\}\vdash C\) and
\(\Theta\cup\{B\}\vdash C\).  DT turns the last two into
\(\Theta\vdash A\to C\) and \(\Theta\vdash B\to C\), and item (iv) gives
\(\Theta\vdash C\).

\emph{\(\to I\)}, discharging occurrences of \(\alpha\) in \(D_1\): by
(\(\ast\)), \(\Gamma_1\subseteq\operatorname{oa}(D)\cup\{\alpha\}\), so the
hypothesis gives \(\Theta\cup\{A\}\vdash B\), and DT gives
\(\Theta\vdash A\to B=\operatorname{tr}_{\mathsf H}(\alpha\to\beta)\).

\emph{\(\to E\)}: the hypotheses give \(A\to B\) and \(A\); one MP step gives
\(B\).

\emph{\(\neg I\)}, discharging occurrences of \(\alpha\) in \(D_1\): by
(\(\ast\)), \(\Gamma_1\subseteq\operatorname{oa}(D)\cup\{\alpha\}\), so the
hypothesis gives \(\bot_{\mathsf H}=\operatorname{tr}_{\mathsf H}(\bottom)\)
from \(\Theta\cup\{A\}\); DT gives \(\Theta\vdash A\to\bot_{\mathsf H}\), and
item (v) gives \(\Theta\vdash\neg A=\operatorname{tr}_{\mathsf H}(\neg\alpha)\).

\emph{\(\neg E\)}: the hypotheses give \(A\) and \(\neg A\); item (vi)
gives \(\bot_{\mathsf H}\).

\emph{\(\bottom E\)}: the hypothesis gives \(\bot_{\mathsf H}\); item (vii)
gives the translation of the conclusion.

\emph{RAA}, discharging occurrences of \(\neg\alpha\) in \(D_1\): by
(\(\ast\)), \(\Gamma_1\subseteq\operatorname{oa}(D)\cup\{\neg\alpha\}\); since
\(\operatorname{tr}_{\mathsf H}(\neg\alpha)=\neg A\), the hypothesis gives
\(\bot_{\mathsf H}\) from \(\Theta\cup\{\neg A\}\); DT gives
\(\Theta\vdash\neg A\to\bot_{\mathsf H}\), and item (viii) gives
\(\Theta\vdash A\).

These are the twelve rules of \texttt{sheet:\allowbreak nd-cpl}
together with the assumption display.  For the consequence, if
\(\operatorname{oa}(D)\subseteq\Gamma\) then
\(\overline{\operatorname{oa}(D)}\subseteq\operatorname{tr}_{\mathsf H}[\Gamma]\),
and \cref{thm:p02-w02b-p02-010-hilbert-structural}(1) applies.
\end{proof}
\end{annexitem}

\begin{annexitem}{annex:p02-w02b-p02-010-hilbert-to-lk}{thm:p02-w02b-p02-010-hilbert-to-lk}{From Hilbert derivations to sequent derivations}
\begin{proof}
By course-of-values induction on \(i\) construct \(\Pi_i\) and the
derivation.  Rule names are those of the inherited sequent calculus; an
identity sequent is \(\Gamma',\varphi\Rightarrow\varphi,\Delta'\) for any
finite multisets \(\Gamma',\Delta'\).

\emph{Premise line}: \(\Pi_i=\theta_i\), and \(\theta_i\Rightarrow\theta_i\) is
an identity sequent.

\emph{A1 line}: \(\Pi_i\) is empty, and
\[
\begin{array}{rll}
1.&\varphi,\psi\Rightarrow\varphi&\text{identity}\\
2.&\varphi\Rightarrow\psi\to\varphi&\to R,1\\
3.&\Rightarrow\varphi\to(\psi\to\varphi)&\to R,2
\end{array}
\]
is a derivation; it needs no case split, since \(\varphi,\varphi\Rightarrow\varphi\)
is also an identity sequent.

\emph{A2 line}, with \(u=\varphi\to(\psi\to\chi)\), \(v=\varphi\to\psi\):
\[
\begin{array}{rll}
1.&\varphi\Rightarrow\chi,\psi,\varphi&\text{identity}\\
2.&\varphi,\psi\Rightarrow\chi,\psi&\text{identity}\\
3.&\varphi,v\Rightarrow\chi,\psi&\to L,1,2\\
4.&v,\varphi,\chi\Rightarrow\chi&\text{identity}\\
5.&v,\varphi,\psi\to\chi\Rightarrow\chi&\to L,3,4\\
6.&v,\varphi\Rightarrow\chi,\varphi&\text{identity}\\
7.&v,\varphi,u\Rightarrow\chi&\to L,6,5\\
8.&u,v\Rightarrow\varphi\to\chi&\to R,7\\
9.&u\Rightarrow v\to(\varphi\to\chi)&\to R,8\\
10.&\Rightarrow u\to(v\to(\varphi\to\chi))&\to R,9
\end{array}
\]
Each \(\to L\) step has the shared context of the rule: at line 3 the
context is \(\varphi\Rightarrow\chi,\psi\) with principal formula \(v\); at
line 5 it is \(v,\varphi\Rightarrow\chi\) with principal \(\psi\to\chi\); at
line 7 it is \(v,\varphi\Rightarrow\chi\) with principal \(u\).

\emph{A3 line}, with \(u=\neg\psi\to\neg\varphi\), \(v=\neg\psi\to\varphi\):
\[
\begin{array}{rll}
1.&u,\psi\Rightarrow\psi&\text{identity}\\
2.&u\Rightarrow\psi,\neg\psi&\neg R,1\\
3.&\varphi,\psi\Rightarrow\psi&\text{identity}\\
4.&\varphi\Rightarrow\psi,\neg\psi&\neg R,3\\
5.&\varphi\Rightarrow\psi,\varphi&\text{identity}\\
6.&\varphi,\neg\varphi\Rightarrow\psi&\neg L,5\\
7.&\varphi,u\Rightarrow\psi&\to L,4,6\\
8.&u,v\Rightarrow\psi&\to L,2,7\\
9.&u\Rightarrow v\to\psi&\to R,8\\
10.&\Rightarrow u\to(v\to\psi)&\to R,9
\end{array}
\]
At line 7 the context is \(\varphi\Rightarrow\psi\) with principal \(u\); at
line 8 it is \(u\Rightarrow\psi\) with principal \(v\).  None of these three
derivations uses Cut.

\emph{MP line}: \(\theta_i\) follows from \(\theta_j\) and
\(\theta_k=\theta_j\to\theta_i\).  The identity sequents
\(\theta_j\Rightarrow\theta_i,\theta_j\) and
\(\theta_j,\theta_i\Rightarrow\theta_i\) give
\(\theta_j,\theta_j\to\theta_i\Rightarrow\theta_i\) by \(\to L\).  Cut with
\(\Pi_k\Rightarrow\theta_j\to\theta_i\) on \(\theta_j\to\theta_i\) gives
\(\Pi_k,\theta_j\Rightarrow\theta_i\), and Cut with \(\Pi_j\Rightarrow\theta_j\)
on \(\theta_j\) gives \(\Pi_j,\Pi_k\Rightarrow\theta_i\).  Put
\(\Pi_i=\Pi_j,\Pi_k\); its support is contained in \(\Gamma\).

Finally, from \(\Pi_n\Rightarrow\theta_n\), apply \(C_L\) to each formula
of multiplicity greater than one until every member of
\(\operatorname{supp}(\Pi_n)\) has multiplicity one, then apply \(W_L\) once
for each member of \(\Gamma_0\setminus\operatorname{supp}(\Pi_n)\).  This
derives \(\Gamma_0\Rightarrow\theta_n\).  For finite \(\Gamma\), take
\(\Gamma_0=\Gamma\).
\end{proof}
\end{annexitem}

\begin{annexitem}{annex:p02-w02b-p02-010-lk-to-hilbert}{thm:p02-w02b-p02-010-lk-to-hilbert}{From sequent derivations to Hilbert derivations}
\begin{proof}
The first half of the proof in \cref{annex:p7.calculi.3} constructs, by
induction on a derivation of \(\Gamma\Rightarrow\Delta\) on
\texttt{sheet:\allowbreak lk}, a natural-deduction derivation of
\(D_\Delta\) whose open assumptions lie in \(\operatorname{supp}(\Gamma)\).
By \cref{def:p02-w02b-p02-010-setting}(3), an
\(\mathsf{LK}_{\mathsf{CPL}}\) derivation contains only identity and
\(\bottom\)-sequents and the rules \(W_L,W_R,C_L,C_R\), Cut, Mix, and the
left and right rules for \(\neg,\land,\lor,\to\).

\emph{Removing Mix.}  Each Mix application is replaced by contractions and
one Cut with the same premises and the same conclusion.  Let the premises be
\(\Gamma\Rightarrow\Delta,\varphi^{m}\) and \(\varphi^{n},\Sigma\Rightarrow\Pi\)
with \(m,n\ge 1\), and the conclusion \(\Gamma,\Sigma\Rightarrow\Delta,\Pi\).
For \(k=m,m-1,\ldots,2\), the \(C_R\) instance with context
\(\Gamma\Rightarrow\Delta,\varphi^{k-2}\) takes
\(\Gamma\Rightarrow\Delta,\varphi^{k}\) to
\(\Gamma\Rightarrow\Delta,\varphi^{k-1}\); after these \(m-1\) steps the
sequent is \(\Gamma\Rightarrow\Delta,\varphi\).  In the same way \(n-1\)
instances of \(C_L\) take \(\varphi^{n},\Sigma\Rightarrow\Pi\) to
\(\varphi,\Sigma\Rightarrow\Pi\).  Cut on \(\varphi\) gives
\(\Gamma,\Sigma\Rightarrow\Delta,\Pi\), the Mix conclusion; occurrences of
\(\varphi\) in \(\Delta\) or \(\Sigma\) are untouched by both the Mix and its
replacement.  The replacement changes no other node, so replacing the Mix
applications one at a time, each step lowering their number by one, yields a
derivation of the same end sequent whose rules are among those listed above
without Mix.  The construction below is applied to this derivation.

For the remaining rules, which are exactly the Mix-free rules of
\(\mathsf{LK}_{\mathsf{CPL}}\), the construction of \cref{annex:p7.calculi.3}
uses the assumption display, the padding detour
(\(\land I\) followed by \(\land E_2\)), finite-disjunction inclusion and
reassociation (by \(\lor I_1\), \(\lor I_2\), \(\lor E\) and, for an empty
smaller support, \(\bottom E\)), and the rules \(\land I\), \(\land E_1\),
\(\land E_2\), \(\lor E\), \(\to I\), \(\to E\), \(\neg I\), \(\neg E\),
\(\bottom E\) and RAA.  The quantifier, alpha and equality rules occur there
only in the cases for \(\forall L\), \(\forall R\), \(\exists L\) and
\(\exists R\), which do not arise.  Every formula introduced by these
propositional cases is built from formulas of the given derivation by
\(\lor\), \(\neg\), \(\land\), \(\to\) and \(\bottom\), so it is
propositional.  The rules used have the same displays and nonvacuous
discharge conditions on \texttt{sheet:\allowbreak nd-cpl} as on
\texttt{sheet:\allowbreak nd-fol}.  Hence the constructed tree
is an \(\mathsf{ND}_{\mathsf{CPL}}\) derivation \(E\) with root \(D_\Delta\)
and \(\operatorname{oa}(E)\subseteq\operatorname{supp}(\Gamma)\).
\Cref{thm:p02-w02b-p02-010-nd-to-hilbert} applied to \(E\) gives
\(\operatorname{tr}_{\mathsf H}[\operatorname{oa}(E)]\derives_{\mathsf H_{\mathsf{CPL}}}
\operatorname{tr}_{\mathsf H}(D_\Delta)\), and
\cref{thm:p02-w02b-p02-010-hilbert-structural}(1) enlarges the premise set
to \(\operatorname{tr}_{\mathsf H}[\operatorname{supp}(\Gamma)]\).  When
\(\operatorname{supp}(\Delta)=\{\alpha\}\), \(D_\Delta\) is \(\alpha\) itself
by \cref{def:p07-sequent-nd-translation}.
\end{proof}
\end{annexitem}

\begin{annexitem}{annex:p02-w02b-p02-010-hilbert-equivalence}{thm:p02-w02b-p02-010-hilbert-equivalence}{Hilbert legs of the derivability equivalence}
\begin{proof}
(1) Left to right is \cref{thm:p02-w02b-p02-010-nd-to-hilbert}.  For right
to left, suppose
\(\operatorname{tr}_{\mathsf H}[\Gamma]\derives_{\mathsf H_{\mathsf{CPL}}}
\operatorname{tr}_{\mathsf H}(\alpha)\).  The premise set lies in the
fragment, so \cref{thm:p02-w02b-p02-010-hilbert-to-nd} gives
\(\operatorname{tr}_{\mathsf H}[\Gamma]\derives_{\mathsf{ND}_{\mathsf{CPL}}}
\operatorname{tr}_{\mathsf H}(\alpha)\).  By
\cref{thm:p02-w02b-p02-010-nd-translation}, every \(\gamma\in\Gamma\)
satisfies \(\{\gamma\}\derives_{\mathsf{ND}_{\mathsf{CPL}}}
\operatorname{tr}_{\mathsf H}(\gamma)\), hence
\(\Gamma\derives_{\mathsf{ND}_{\mathsf{CPL}}}\operatorname{tr}_{\mathsf H}(\gamma)\).
\Cref{thm:p02-w02b-p02-010-nd-structural}(3) with empty first set,
\(\Sigma=\operatorname{tr}_{\mathsf H}[\Gamma]\) and \(\Delta=\Gamma\) gives
\(\Gamma\derives_{\mathsf{ND}_{\mathsf{CPL}}}\operatorname{tr}_{\mathsf H}(\alpha)\).
The second half of \cref{thm:p02-w02b-p02-010-nd-translation} gives
\(\{\operatorname{tr}_{\mathsf H}(\alpha)\}\derives_{\mathsf{ND}_{\mathsf{CPL}}}\alpha\),
and composition with \(\Sigma=\{\operatorname{tr}_{\mathsf H}(\alpha)\}\)
and \(\Delta=\Gamma\) gives \(\Gamma\derives_{\mathsf{ND}_{\mathsf{CPL}}}\alpha\).

(2) Let \(\Gamma\) be finite, read as a multiset of multiplicity one, so
\(\operatorname{supp}(\Gamma)=\Gamma\).  Left to right is
\cref{thm:p02-w02b-p02-010-lk-to-hilbert} with \(\Delta=\alpha\).  For
right to left, part (1) gives \(\Gamma\derives_{\mathsf{ND}_{\mathsf{CPL}}}\alpha\),
witnessed by a derivation \(E\) with \(\operatorname{oa}(E)\subseteq\Gamma\).
The second half of the proof in \cref{annex:p7.calculi.3} turns \(E\) into a
derivation of \(\operatorname{oa}(E)\Rightarrow\alpha\) on
\texttt{sheet:\allowbreak lk}.  For the rules of
\(\mathsf{ND}_{\mathsf{CPL}}\) it uses only contexted identity, the
\(\bottom\)-sequent, weakening, contraction, Cut, and the left and right
rules for \(\neg,\land,\lor,\to\), applied to formulas of \(E\) and
\(\bottom\); so the result is an \(\mathsf{LK}_{\mathsf{CPL}}\) derivation.
Left weakening by the members of \(\Gamma\setminus\operatorname{oa}(E)\)
gives \(\Gamma\Rightarrow\alpha\).

(3) \(\operatorname{tr}_{\mathsf H}\) is the identity on
\(\mathsf{Fm}_{\{\neg,\to\}}\), by structural induction using its clauses for
variables, \(\neg\) and \(\to\).  So for
\(\Gamma\cup\{\alpha\}\subseteq\mathsf{Fm}_{\{\neg,\to\}}\), parts (1) and (2)
read \(\Gamma\derives_{\mathsf{ND}_{\mathsf{CPL}}}\alpha\) iff
\(\Gamma\derives_{\mathsf H_{\mathsf{CPL}}}\alpha\), and, for finite
\(\Gamma\), \(\Gamma\Rightarrow\alpha\) is derivable iff
\(\Gamma\derives_{\mathsf H_{\mathsf{CPL}}}\alpha\).

No step uses a semantic notion, so the corollary does not depend on
soundness or completeness of any of the three calculi.
\end{proof}
\end{annexitem}
